\chapter{Conclusão}\label{cap:conclusao}
\glsresetall

% =====================================================================
% Estrutura (espelha o modelo institucional): recapitulacao + Contribuicoes +
% Limitacoes + Trabalhos Futuros. Como o estudo esta em estruturacao, a
% conclusao e apresentada em termos ESPERADOS; os trechos a fixar apos a coleta
% estao como PLACEHOLDER.
% =====================================================================

Este trabalho apresentou o desenvolvimento do \textit{Cozy Pong}, um jogo de \gls{rv} que une tênis de mesa e ritmo sob músicas \textit{lo-fi}, e propôs um protocolo experimental para caracterizar seu efeito sobre a regulação do estresse e da ansiedade. A jogabilidade central, o sistema de mapas rítmicos, a pontuação e a retroalimentação multissensorial foram implementados na \textit{engine} Unity para o Meta Quest 3S, e sobre essa base foram definidas quatro condições experimentais: o jogo em configuração relaxante, o mesmo jogo em configuração animada, a audição da mesma trilha sem jogar e uma sessão de meditação guiada no mesmo dispositivo. O protocolo adota desenho intra-sujeito randomizado e contrabalanceado por quadrado de Williams, com indução de estresse padronizada, bateria psicométrica integralmente composta por escalas de uso livre validadas em português brasileiro e aferição psicofisiológica por \gls{vfc} e \gls{eda}.

\ph{enunciar a conclusão principal do estudo após a coleta e a análise, respondendo diretamente às perguntas de pesquisa PP1, PP2 e PP3 e às hipóteses H1 a H3.}

% =====================================================================
\section{Contribuições}\label{sec:conc-contribuicoes}

As contribuições esperadas organizam-se em quatro frentes. A primeira é o próprio artefato, o \textit{Cozy Pong}, um \textit{exergame} de \gls{rv} rítmico com \textit{beatmaps} editáveis e retroalimentação multissensorial, acompanhado das duas configurações de jogabilidade e da cena de meditação implementadas no mesmo aplicativo, o que constitui, por si, um instrumento experimental reutilizável.

A segunda é o desenho experimental, construído para que cada contraste varie um único bloco de fatores: a trilha musical é mantida constante no contraste que isola o jogo, o dispositivo e o cenário são mantidos constantes no contraste que isola a caracterização relaxante, e o meio de entrega é mantido constante no contraste com a meditação guiada. Essa estrutura permite atribuir eventuais diferenças a fatores identificados, em vez de apenas constatar que a exposição ao jogo é seguida de melhora.

A terceira é metodológica e diz respeito à seleção de instrumentos. O trabalho documenta a substituição do instrumento predominante na literatura, indisponível para uso livre e de aplicação privativa de psicólogos no Brasil, por uma bateria integralmente livre e validada em português brasileiro, contribuição diretamente aproveitável por outros estudos brasileiros de computação e saúde que enfrentem a mesma restrição.

A quarta é a evidência psicofisiológica esperada, que quantifica como um \textit{exergame} rítmico se posiciona frente à prática de relaxamento digital hoje difundida. \ph{ajustar a redação das contribuições à luz dos resultados obtidos.}

% =====================================================================
\section{Limitações}\label{sec:conc-limitacoes}

\begin{itemize}
    \item O estudo de exposição única em participantes saudáveis não autoriza alegações de eficácia terapêutica, mas sim estimativas do efeito de uma intervenção breve para a regulação de estresse.
    \item O contraste com a audição passiva separa o efeito do jogo do efeito da música, mas os componentes internos do jogo, a imersão, a estrutura rítmica, o movimento e o engajamento lúdico, variam em bloco, de modo que uma eventual vantagem é atribuível ao jogo como um todo e não à sincronização rítmica isoladamente.
    \item Nem a audição da trilha nem a meditação guiada são controles neutros: ambas possuem efeito relaxante próprio e funcionam como comparadores ativos. O teste é, portanto, conservador, e a ausência de diferença não indica ineficácia da intervenção, apenas que ela não supera essas alternativas.
    \item O esforço físico não é comparável entre as condições, de modo que parte da resposta fisiológica pode refletir intensidade de atividade e não regulação emocional. O uso de janelas de recuperação e do esforço percebido como covariável atenua, mas não elimina, essa ameaça.
    \item As condições não admitem mascaramento, pois o participante identifica imediatamente qual atividade realiza, o que expõe o estudo a características de demanda e a efeitos de expectativa. Some-se a atenuação esperada da resposta ao estressor ao longo das quatro sessões, estimada pelo modelo mas não anulada.
    \item A amostra de conveniência de jovens adultos e universitários limita a generalização, e a experiência prévia com meditação pode moderar a resposta à condição de comparação.
    \item A exigência de quatro sessões por participante impõe carga elevada e constitui a principal ameaça à retenção amostral e à viabilidade do estudo.
    % PLACEHOLDER: acrescentar limitacoes efetivamente observadas apos a coleta
    % (p.ex.\ taxa de falha de sensor, atrito amostral, cybersickness).
\end{itemize}

% =====================================================================
\section{Trabalhos Futuros}\label{sec:conc-futuros}

\begin{itemize}
    \item Acrescentar uma variante do próprio jogo com mapeamento rítmico dessincronizado, de modo a decompor a contribuição da sincronização rítmica, que o desenho atual agrega ao restante do jogo.
    \item Comparar a meditação guiada em \gls{rv} com a meditação presencial, contraste deliberadamente excluído deste desenho por confundir técnica e meio de entrega em um único termo.
    \item Estender a avaliação a populações com estresse ou ansiedade clinicamente relevantes, com acompanhamento e comparadores apropriados, o que exigiria trilhas éticas adicionais.
    \item Investigar intervenções repetidas e efeitos de médio prazo, para além da exposição única, cenário no qual a fruição superior de um jogo pode traduzir-se em maior adesão do que a de técnicas convencionais.
    \item Explorar a personalização da trilha \textit{lo-fi} e do mapeamento rítmico às preferências do participante, dado que referências musicais distintas podem produzir respostas distintas.
    % PLACEHOLDER: acrescentar direcoes sugeridas pelos resultados.
\end{itemize}
